%=========================================================================
\chapter{Building, Installing and Using WIR}
\label{chap:build}
\renewcommand{\sectiontitle}{}
%=========================================================================

Once the build system has been set up as described in \cref{chap:bootstrap} before, building \gls{WIR} is as easy as depicted in \cref{lst:buildmake} of \cref{lst:build} -- a simple call of \texttt{make} in the \texttt{BUILD} directory is sufficient.

\begin{lstlisting}[style=commandline,float=h,caption={WIR Build using \texttt{make}},label=lst:build]
@/tmp/WIR/BUILD@> make -j 8Ü\label{lst:buildmake}Ü
make[1]: Entering directory '/tmp/WIR/BUILD'
Making all in LIBUSEFUL
make[2]: Entering directory '/tmp/WIR/BUILD/LIBUSEFUL'
make  all-recursive
make[3]: Entering directory '/tmp/WIR/BUILD/LIBUSEFUL'
Making all in libuseful
make[4]: Entering directory '/tmp/WIR/BUILD/LIBUSEFUL/libuseful'
libtool: compile:  g++ -DHAVE_CONFIG_H -I. -I../../../LIBUSEFUL/libuseful -I.. -Og -g -Wall -Wextra -std=c++17 -m64 -MT auxfuncs.lo -MD -MP -MF .deps/auxfuncs.Tpo -c ../../../LIBUSEFUL/libuseful/auxfuncs.cc  -fPIC -DPIC -o .libs/auxfuncs.o
libtool: compile:  g++ -DHAVE_CONFIG_H -I. -I../../../LIBUSEFUL/libuseful -I.. -Og -g -Wall -Wextra -std=c++17 -m64 -MT debugmacros.lo -MD -MP -MF .deps/debugmacros.Tpo -c ../../../LIBUSEFUL/libuseful/debugmacros.cc  -fPIC -DPIC -o .libs/debugmacros.o
libtool: compile:  g++ -DHAVE_CONFIG_H -I. -I../../../LIBUSEFUL/libuseful -I.. -Og -g -Wall -Wextra -std=c++17 -m64 -MT exceptions.lo -MD -MP -MF .deps/exceptions.Tpo -c ../../../LIBUSEFUL/libuseful/exceptions.cc  -fPIC -DPIC -o .libs/exceptions.o
[...]
libtool: link: ( cd ".libs" && rm -f "libwir.la" && ln -s "../libwir.la" "libwir.la" )
make[2]: Leaving directory '/tmp/WIR/BUILD'
Making all in doc
make[2]: Entering directory '/tmp/WIR/BUILD/doc'
make[2]: Leaving directory '/tmp/WIR/BUILD/doc'
make[1]: Leaving directory '/tmp/WIR/BUILD'
\end{lstlisting}

The example from \cref{lst:buildmake} of \cref{lst:build} uses 8 processor cores of the host system in parallel, which significantly speeds up the build procedure.
After \gls{WIR} has been built, its unit tests can be executed (cf. \cref{lst:testsmake} of \cref{lst:tests}).
These tests are designed to cover true-positive and true-negative scenarios which are represented by \texttt{PASS} and \texttt{XFAIL} in \cref{lst:tests}.
When running the unit tests, obviously no false-positives (\texttt{XPASS}) or false-negatives (\texttt{FAIL}) should occur.
If configured accordingly (cf. \cref{tab:configure}), the HTML reference documentation generated by \texttt{doxygen} can be inspected by opening file \texttt{WIR/BUILD/doc/html/index.html} in a browser.

After the successful build of \gls{WIR}, the complete software package can be installed in the directory configured as installation location with the \texttt{-{}-prefix} option of \texttt{configure} (cf. \cref{tab:configure}):

\begin{lstlisting}[style=commandline,float=h]
@/tmp/WIR/BUILD@> make install
\end{lstlisting}

This creates a directory hierarchy as depicted in \cref{fig:wirinstalldirectory}.
The \texttt{etc} directory contains sub-directories for each \gls{ISA}, and the system-level configuration files are installed therein.
The \texttt{include} directory tree contains all C++ header files of \gls{WIR}, following the same structure as shown in \cref{fig:wircodedirectory}.
The built binary libraries \texttt{libuseful.so} and \texttt{libwir.so} (and/or their static counterparts if configured accordingly) reside in sub-directory \texttt{lib} after installation.
Finally, sub-directory \texttt{share} contains the authors, license and version information of the software package.

\begin{lstlisting}[style=commandline,float=t,caption={WIR Unit Tests},label=lst:tests]
@/tmp/WIR/BUILD@> make -j 8 checkÜ\label{lst:testsmake}Ü
Making check in LIBUSEFUL
[...]
Making check in tests
make[1]: Entering directory '/tmp/WIR/BUILD/tests'
[...]
^PASS^: tc1796MemoryLayout2.sh
^PASS^: tcFunction3.sh
^PASS^: tcCompilationUnit1.sh
^PASS^: tcOperations2.sh
^PASS^: tcBasicBlock1.sh
^============================================================================^
^Testsuite summary for wir 2.0^
^============================================================================^
~~# TOTAL: 126~~
^# PASS:  89^
~~# SKIP:  0~~
§# XFAIL: 37§
~~# FAIL:  0~~
~~# XPASS: 0~~
~~# ERROR: 0~~
^============================================================================^
make[5]: Leaving directory '/tmp/WIR/BUILD/arÜÜch/tricore/tests'
make[2]: Leaving directory '/tmp/WIR/BUILD/arÜÜch/tricore'
make[1]: Leaving directory '/tmp/WIR/BUILD/arÜÜch'
make[1]: Leaving directory '/tmp/WIR/BUILD'
\end{lstlisting}

\begin{figure}[!ht]
  \begin{center}
    \begin{forest}
      pic dir tree,
      pic root,
      for tree={% folder icons by default; override using file for file icons
        directory,
      },
      [\emph{<INSTALLDIR>}
        [etc
          [arm]
          [\dots]
          [tricore
            [tc1796.sys, file]
            [\dots, file]
          ]
        ]
        [include
          [analyses]
          [\dots]
          [wir
            [wir.h, file]
            [\dots, file]
          ]
        ]
        [lib
          [libuseful.so, file]
          [libwir.so, file]
        ]
        [share
          [libuseful]
          [wir
            [AUTHORS, file]
            [LICENSE, file]
            [\dots, file]
          ]
        ]
      ]
    \end{forest}
  \end{center}
  \caption{Directory Structure of Installed WIR Library}
  \label{fig:wirinstalldirectory}
\end{figure}

This installation can now be used by \gls{WIR}-based applications.
\cref{lst:testprogram} shows a very simple test program that includes the generic and \gls{ISA}-independent headers in \cref{lst:testprograminclude}.
All C++ classes and types of the \gls{WIR} library are placed into a dedicated C++ namespace called \texttt{WIR}.
\cref{lst:testprogramnamespace} enables the global usage of this namespace so that no explicit qualifiers need to be given in the subsequent code.\footnote{In the rest of this entire document, we from now on assume in all code listings that this namespace \texttt{WIR} is always used. All following listings will thus not contain any qualifier \texttt{WIR::} in order to improve code legibility.}
Before any use of \gls{WIR} data structures in application code, the library must be initialized (cf. \cref{lst:testprograminit}).
After that, the test program creates an instance of a \texttt{WIR\_System} based on the MIPS SPIM generic \gls{ISA} model installed in file \texttt{\emph{<INSTALLDIR>}/etc/generic/genericmips.sys} (cf. \cref{lst:testprogramsys}).

The test program can finally be compiled and linked against the installed \gls{WIR} library as follows:

\begin{lstlisting}[style=commandline,float=h]
@/tmp@> export LD_LIBRARY_PATH=<<<INSTALLDIR>>>/lib:$LD_LIBRARY_PATHÜ\label{lst:usewir1}Ü

@/tmp@> g++ wirtest.cc -o wirtest.bin -I<<<INSTALLDIR>>>/include -L<<<INSTALLDIR>>>/lib -luseful -lwirÜ\label{lst:usewir2}Ü
\end{lstlisting}

Setting the variable \texttt{LD\_LIBRARY\_PATH} as shown in \cref{lst:usewir1} above (again assuming a Bourne-compatible shell) is required so that the host system's dynamic linker is able to find the installed binary libraries.
\cref{lst:usewir2} afterwards invokes the \texttt{g++} compiler with the proper include and library arguments.

\begin{lstlisting}[style=code,float=b,caption={Simple WIR Test Program \texttt{wirtest.cc}},label=lst:testprogram]
// Include WIR headers.
^#include^ §<wir/wir.h>§Ü\label{lst:testprograminclude}Ü

using namespace WIR;Ü\label{lst:testprogramnamespace}Ü

@int@ main@( void )@
{
  // Initialize WIR.
  WIR_Init@()@;Ü\label{lst:testprograminit}Ü

  // Create test WIR system.
  WIR_TaskManager t;
  WIR_System sys@(@ "genericmips.sys", t @)@;Ü\label{lst:testprogramsys}Ü

  sys.setName@(@ "Test WIR system" @)@;

  return@(@ §0§ @)@;
}
\end{lstlisting}

\cleardoubleoddpage
